
par (i) @@M0@@ (iv) prenons une surjection @@M1@@. Quitte à couper @@M2@@ en morceaux, de sorte que @@M3@@ devienne plat sur @@M4@@, la formation de la suite exacte @@M5@@ commute au passage aux fibres. Par (i) @@M6@@ (iii) et (i) @@M7@@ (iv), prenons @@M8@@ assez grand pour que tous les @@M9@@ soient nuls et qu'il existe des surjections @@M10@@ et @@M11@@. Alors les applications @@M12@@ sont surjectives. Soit @@M13@@ un homomorphisme et soit @@M14@@ la composition @@M15@@; alors @@M16@@ provient d'un homomorphisme @@M17@@, et @@M18@@ et @@M19@@ ont le même conoyau. Par conséquent, @@M20@@ est contenue dans la famille des classes des conoyaux d'homomorphismes de la forme @@M21@@.

Toujours quitte à couper @@M22@@ en morceaux, on peut supposer en plus que @@M23@@ (donc @@M24@@) et tous les @@M25@@ (où @@M26@@ est le morphisme structural) sont plats sur @@M27@@, afin que la formation de @@M28@@ commute au changement de base quelconque @@M29@@ (EGA III 6.9.9.2). Soit @@M30@@, alors @@M31@@ porte canoniquement une suite exacte ``universelle''
@@M32@@
et @@M33@@ limite @@M34@@.

\subsection*{2. Plusieurs lemmes techniques}
\textbf{2.1.} Soient @@M35@@ un corps et @@M36@@ un @@M37@@-schéma propre. Soient @@M38@@ des faisceaux inversibles sur @@M39@@, @@M40@@ la classe de Chern de @@M41@@ et @@M42@@ un sous-schéma fermé de @@M43@@ de dimension @@M44@@. Rappelons qu'on désigne par
